\subsection{Transpos}

Operasi penting lainnya pada matriks adalah mengambil \textbf{transpos}. Untuk matriks $A$, kita menyatakan
\textbf{transpos} dari $A$ dengan $A^T$. Sebelum mendefinisikan transpos secara formal, kita menelaah
operasi ini pada matriks berikut.
\begin{equation*}
\leftB
\begin{array}{cc}
1 & 4 \\
3 & 1 \\
2 & 6
\end{array}
\rightB ^{T}=\allowbreak \leftB
\begin{array}{ccc}
1 & 3 & 2 \\
4 & 1 & 6
\end{array}
\rightB
\end{equation*}

Apa yang terjadi? Kolom pertama menjadi baris pertama dan kolom kedua
menjadi baris kedua. Dengan demikian, matriks $3\times 2$ menjadi matriks $2\times 3$.
Bilangan $4$ semula berada pada baris pertama dan kolom kedua, lalu
berpindah ke baris kedua dan kolom pertama. 

Definisi transpos adalah sebagai berikut.

\begin{definition}{Transpos Suatu Matriks}\label{def:matrixtranspose}
Misalkan $A$ adalah matriks $m\times n$. Maka $A^{T}$, yaitu \textbf{transpos} \index{matriks!transpos} dari $A$, menyatakan matriks $n\times m$
yang diberikan oleh
\begin{equation*}
A^{T} = \leftB a _{ij}\rightB ^{T}= \leftB a_{ji} \rightB
\end{equation*}
\end{definition}

Entri-$\left( i, j  \right)$ dari $A$ menjadi
entri-$\left( j,i \right)$ dari $A^T$. 

Perhatikan contoh berikut.

\begin{example}{Transpos Suatu Matriks}\label{exa:transposematrix}
Hitung $A^T$ untuk matriks berikut
\begin{equation*}
A = \leftB
\begin{array}{rrr}
1 & 2 & -6 \\
3 & 5 & 4
\end{array}
\rightB
\end{equation*}
\end{example}

\begin{solution}
Berdasarkan Definisi \ref{def:matrixtranspose}, kita mengetahui bahwa untuk $A = \leftB a_{ij} \rightB$,
$A^T = \leftB a_{ji} \rightB$. Dengan kata lain, kita menukar posisi baris dan kolom
setiap entri. Entri-$\left( 1, 2 \right)$ menjadi entri-$\left( 2,1 \right)$.

Jadi,
\begin{equation*}
A^T = 
 \leftB
\begin{array}{rr}
1 & 3 \\
2 & 5 \\
-6 & 4
\end{array}
\rightB 
\end{equation*}

Perhatikan bahwa $A$ adalah matriks $2 \times 3$, sedangkan $A^T$ adalah matriks $3 \times 2$.
\end{solution}

Transpos suatu matriks memiliki sifat-sifat penting berikut \index{matriks!sifat transpos}.

\begin{lemma}{Sifat-Sifat Transpos Suatu Matriks}\label{lem:transposeproperties}
Misalkan $A$ adalah matriks $m\times n$, $B$ adalah matriks $n\times p$, serta $r$ dan $s$ adalah skalar. Maka
\begin{enumerate}
\item
$\left(A^{T}\right)^{T} = A$
\item
$\left( AB\right) ^{T}=B^{T}A^{T} $ \label{matrixtranspose1}
\item
$\left( rA+ sB\right) ^{T}=rA^{T}+ sB^{T}$  \label{matrixtranspose2}
\end{enumerate}
\end{lemma}

\ifdefined\showproofs
\begin{proof}
Pertama, kita membuktikan \ref{matrixtranspose1}. Dari Definisi \ref{def:matrixtranspose},
\begin{equation*}
\left(AB\right)^{T} = \leftB (AB) _{ij} \rightB ^{T}=\leftB (AB)_{ji} \rightB=\sum_{k}a_{jk}b_{ki}= \sum_{k}b_{ki}a_{jk} 
\end{equation*}
\begin{equation*}
= \sum_{k}\leftB b_{ik}\rightB^{T}\leftB
a_{kj}\rightB^{T}=\leftB b_{ij}\rightB ^{T} \leftB a_{ij}\rightB^{T} = B^{T}A^{T} 
\end{equation*}
Bukti Rumus \ref{matrixtranspose2} diserahkan sebagai latihan. 
\end{proof}
\fi
Transpos suatu matriks berkaitan dengan topik-topik penting lainnya. Perhatikan definisi berikut.  

\begin{definition}{Matriks Simetris dan Simetris-Miring}\label{def:symmetricandskewsymmetric}
Matriks $n\times n$, yaitu $A$, disebut
\index{matriks!simetris}\textbf{simetris} jika $A=A^{T}.$ Matriks tersebut disebut
\index{matriks!simetris-miring} \textbf{simetris-miring} jika $A=-A^{T}.$
\end{definition}

Kita akan menelaah definisi-definisi ini dalam contoh-contoh berikut.

\begin{example}{Matriks Simetris}\label{exa:symmetricmatrix}
Misalkan
\begin{equation*}
A=\leftB
\begin{array}{rrr}
2 & 1 & 3 \\
1 & 5 & -3 \\
3 & -3 & 7
\end{array}
\rightB 
\end{equation*}
Gunakan Definisi \ref{def:symmetricandskewsymmetric} untuk menunjukkan bahwa $A$ simetris. 
\end{example}

\begin{solution}
Berdasarkan Definisi \ref{def:symmetricandskewsymmetric}, kita perlu menunjukkan bahwa $A = A^T$.
Sekarang, dengan menggunakan Definisi \ref{def:matrixtranspose},
\begin{equation*}
A^{T} = \leftB
\begin{array}{rrr}
2 & 1 & 3 \\
1 & 5 & -3 \\
3 & -3 & 7
\end{array}
\rightB
\end{equation*}
Jadi, $A = A^{T}$, sehingga $A$ simetris.
\end{solution}

\begin{example}{Matriks Simetris-Miring}\label{exa:skewsymmetricmatrix}
Misalkan
\begin{equation*}
A=\leftB
\begin{array}{rrr}
0 & 1 & 3 \\
-1 & 0 & 2 \\
-3 & -2 & 0
\end{array}
\rightB 
\end{equation*}
Tunjukkan bahwa $A$ simetris-miring.
\end{example}

\begin{solution} Berdasarkan Definisi \ref{def:symmetricandskewsymmetric},
\begin{equation*}
A^{T} = \leftB
\begin{array}{rrr}
0 & -1 & -3\\
1 &  0 & -2\\
3 &  2 &  0
\end{array}
\rightB 
\end{equation*}

Anda dapat melihat bahwa setiap entri $A^T$ sama dengan $-1$ kali entri yang sama dari $A$.
Jadi, $A^{T} = - A$ dan berdasarkan Definisi \ref{def:symmetricandskewsymmetric}, $A$ simetris-miring.
\end{solution}
